\documentclass[conference]{IEEEtran}

\usepackage{cite}
\usepackage{amsmath,amssymb}
\usepackage{graphicx}
\usepackage{booktabs}
\usepackage{array}
\usepackage{url}
\usepackage{float}
\usepackage[section]{placeins}
\usepackage{balance}
\usepackage{microtype}

\renewcommand{\ttdefault}{cmtt}
\graphicspath{{figures/}}

\newcommand{\paperfigure}[4]{%
\begin{figure}[H]
\centering
\IfFileExists{figures/#1}{%
  \includegraphics[width=#2\columnwidth]{figures/#1}%
}{%
  \fbox{\parbox[c][1.38in][c]{0.92\columnwidth}{%
    \centering\small
    \textbf{Figure unavailable in this build.}\\[4pt]
    #3
  }}%
}
\caption{#3}
\label{#4}
\end{figure}
}

\title{Real-Time SCADA Intrusion Detection Under Communication Constraints: Latency, Payload, and Generalization in Industrial Edge Environments}

\author{
\IEEEauthorblockN{Digdarshan Subedi}
\IEEEauthorblockA{
Department of Computer Science\\
University of Wisconsin--Whitewater\\
Whitewater, WI, USA\\
SubediD30@uww.edu}
\and
\IEEEauthorblockN{Nikesh Tiwari}
\IEEEauthorblockA{
School of Business \& Technology\\
Webster University\\
St. Louis, MO, USA\\
nikeshtiwari@webster.edu}
\and
\IEEEauthorblockN{Bibek Ale}
\IEEEauthorblockA{
School of Business \& Technology\\
Webster University\\
St. Louis, MO, USA\\
Bibekale@webster.edu}
}

\begin{document}
\maketitle

\begin{abstract}
Digital substations increasingly rely on continuous communication among
sensing, protection, and supervisory components to support timely operational
decisions. Intrusion detectors within this workflow must respond quickly,
operate with the telemetry visible at their processing location, and remain
reliable under conditions not observed during development. However, SCADA
intrusion detection is commonly evaluated as an offline classification task,
separating model accuracy from communication payload, processing latency,
telemetry availability, and cross-run transfer.

We examine this mismatch through a communication-aware evaluation of four
lightweight detectors on the MSU/ORNL Power System Attack Dataset. Using
fifteen-fold leave-one-run-out testing, we study unseen-run discrimination,
batch-size-one pipeline latency, modeled payload, missing telemetry, and
edge-visible versus central feature conditions. Random Forest ROC-AUC decreases
from 0.972 under random splitting to 0.677 under run-aware evaluation.
Logistic Regression, XGBoost, and a compact multilayer perceptron meet the
10~ms and 16.7~ms P95 study references, whereas Random Forest meets only the
100~ms reference. A training-fold-selected 16-feature condition reduces
modeled payload from 512 to 64 bytes while preserving comparable or higher
mean ROC-AUC for three models. Central telemetry improves clean-condition
discrimination, but relay-associated feature loss narrows or removes this
advantage.

These findings show that low-latency inference alone does not establish the
usefulness of a real-time industrial detector. Credible evaluation must
consider unseen-run transfer, telemetry efficiency, placement-dependent
visibility, degradation resilience, and false-alarm burden as connected
properties of the communication workflow.
\end{abstract}


\vspace{0.5em}

\begin{IEEEkeywords}
real-time industrial communications, SCADA security, intrusion detection,
edge computing, industrial IoT, telemetry constraints, digital substations,
machine learning
\end{IEEEkeywords}


\section{Introduction}

Digital substations depend on time-sensitive communication among phasor
measurement units (PMUs), protective relays, control systems, and security
monitors. An intrusion detector in this workflow must not only classify
telemetry accurately, but also operate on the measurements available at its
location, meet practical processing budgets, and avoid excessive false alarms.

These requirements create a real-time communication trade-off. Edge detection
can reduce delay and communication dependence, but local models may have
limited visibility. Central detection offers broader telemetry, but depends on
timely delivery and may degrade under missing measurements or relay-channel
loss. Reducing features can lower payload and processing cost, but may also
remove useful attack evidence.

SCADA intrusion detection is nevertheless often evaluated as an offline
classification task. Random row-level splitting can place records from the
same collection run in both training and test sets, overstating performance
under unseen conditions. Latency, payload, telemetry loss, and detector
placement are also commonly evaluated separately from detection quality.

We therefore frame SCADA intrusion detection as a
\emph{communication-budgeted service}. Using fifteen-fold leave-one-run-out
testing, we evaluate Logistic Regression, Random Forest, XGBoost, and a compact
multilayer perceptron across unseen-run discrimination, batch-size-one latency,
modeled payload, telemetry loss, and edge-visible versus central feature
conditions. These placement conditions are analytical proxies rather than a
deployed substation architecture.

\subsection{Research Questions}

\begin{itemize}
\item \textbf{RQ1:} Which detector pipelines meet the study's latency
references?

\item \textbf{RQ2:} How does feature reduction affect payload, latency, and
unseen-run discrimination?

\item \textbf{RQ3:} How does telemetry loss affect detector performance?

\item \textbf{RQ4:} When does central telemetry outperform edge-visible
telemetry?
\end{itemize}

\subsection{Contributions}

This paper contributes:

\begin{itemize}
\item A communication-aware evaluation of SCADA intrusion detection across
latency, payload, telemetry availability, and placement.

\item A leakage-controlled, run-aware protocol in which all preprocessing,
feature selection, training, and threshold selection exclude the held-out run.

\item Evidence that reduced telemetry can preserve competitive unseen-run
performance, while low latency and central visibility alone do not guarantee
robust detection.
\end{itemize}



\section{Related Work}

\subsection{SCADA Intrusion Detection Beyond Classification Accuracy}

Power-system intrusion detection has been studied using physical telemetry,
cyber logs, and hybrid feature sets. The MSU/ORNL benchmark has supported
disturbance and cyber-attack discrimination with synchronized measurements
\cite{borges2014machine,pan2015hybrid,pan2015disturbances}. Later work has
examined ensemble learning, feature selection, deep models, federated
learning, and lightweight detection
\cite{upadhyay2021intrusion,yang2019deep,wang2022lightweight,shao2024automated}.
Surveys further show that dataset design, class balance, and evaluation
protocols strongly shape reported performance
\cite{sahani2023machine,alimi2021review}. Rather than proposing a new
classifier, we study whether lightweight detectors remain useful on unseen
runs and under constrained telemetry.

\subsection{Communication Constraints and Evaluation Design}

IEC~61850 defines communication requirements for power-utility automation,
while IEEE~C37.118.2 specifies synchrophasor data transfer
\cite{iec61850-5,iec61850-90-4,ieee-c37118}. These standards establish the
real-time communication context of this work, although our measurements do not
demonstrate protocol compliance or protection-loop suitability. Low
computational cost alone is insufficient when evaluation mixes collection
conditions or assumes complete telemetry. Arp \emph{et al.} identify leakage
and unrealistic assumptions as recurring weaknesses in machine-learning
security research \cite{arp2022dodo}. We therefore evaluate run-aware
generalization together with latency, payload, telemetry degradation, and
detector placement.



\FloatBarrier
\section{Communication-Budgeted Detection Framework}

Our framework follows the operational path shown in
Fig.~\ref{fig:architecture}: telemetry is observed under processing, payload,
and availability constraints before classification at an edge-visible or
central location. It is guided by three goals. \textbf{G1: credible transfer}
requires testing on a complete unseen collection run. \textbf{G2: visible
resource trade-offs} connects feature count to payload and measured pipeline
latency. \textbf{G3: explicit degradation} evaluates what remains useful when
measurements or relay-associated channels are unavailable.

\paperfigure{fig01_system_architecture.png}{1.00}
{Overview of the communication-budgeted evaluation framework. Digital-substation
telemetry is examined under processing, payload, and availability constraints
before classification with edge-visible R1 measurements or centrally
aggregated features. The figure describes an analytical evaluation workflow,
not a deployed production architecture.}
{fig:architecture}

\FloatBarrier
\subsection{Dataset and Visibility Conditions}

We use the binary-labeled MSU/ORNL power-system attack dataset from the public
Industrial Control System Cyber Attack Datasets collection
\cite{ics-datasets,borges2014machine,pan2015hybrid}. The audited corpus contains fifteen runs and 78{,}377 records: 55{,}663
attack and 22{,}714 normal/natural records. Its 128 usable features comprise
116 PMU/relay measurements and 12 control-panel, Snort, or relay-log
variables. Each file receives a \texttt{run\_id} used only for grouped
evaluation; it is never a model input.

The edge-visible condition uses the 29 R1 variables whose names begin
\texttt{R1-} or \texttt{R1:}; models are retrained on these measurements only.
The central condition uses all 128 features. The R1 result is one predefined
proxy, not an average or best case across relays. Because the public files do
not preserve a valid continuous event sequence, we evaluate record-level
classification and do not claim attack-onset delay, early warning, or
time-to-detection.

\FloatBarrier
\section{Method}

\subsection{Run-Aware Evaluation and Models}

Each leave-one-run-out fold trains on fourteen complete runs and tests on the
remaining run. Imputation, scaling, feature ranking, model fitting, and
threshold selection use training runs only; availability degradation is
applied only to held-out records. Results are unweighted fold means and
standard deviations unless noted. A stratified 70/30 random split with seed 42
is retained only as a contrast baseline.

\paperfigure{fig02_random_vs_runaware.png}{1.00}
{Random row-level splitting versus leave-one-run-out evaluation for comparable
baseline models. The figure highlights how discrimination changes when the
test set is a complete unseen collection run.}
{fig:random-runaware}

All primary models use training-fitted median imputation. Logistic Regression
and the compact MLP also use training-fitted standardization. Table
\ref{tab:model-config} summarizes the configurations; all use seed 42.

\begin{table}[H]
\caption{Primary model configurations.}
\label{tab:model-config}
\centering
\scriptsize
\begin{tabular}{p{0.22\columnwidth}p{0.64\columnwidth}}
\toprule
Model & Main configuration \\
\midrule
Logistic Regression & Standardized; \texttt{lbfgs}, $\ell_2$, $C=1$, 2{,}000
maximum iterations. \\
Random Forest & 200 Gini trees, square-root feature rule, bootstrap enabled,
unrestricted depth. \\
XGBoost & 200 estimators, depth 6, learning rate 0.3, histogram tree method,
four CPU threads. \\
Compact MLP & Standardized; hidden layers $(32,16)$, ReLU, Adam, early
stopping, 200 maximum iterations. \\
\bottomrule
\end{tabular}
\end{table}

\FloatBarrier
\subsection{Processing and Payload Measurements}

Batch-size-one latency is measured on cached fold-01 pipelines using 500
warm-up and 5{,}000 timed calls with \texttt{time.perf\_counter\_ns()}. Timing
includes the fitted preprocessor and estimator \texttt{predict\_proba()} call,
but excludes acquisition, file I/O, serialization, transport, queueing,
visualization, and operator response. Measurements were collected on an ARM64
macOS host with ten cores and 32~GB memory using Python~3.9.6,
scikit-learn~1.6.1, and XGBoost~2.1.4. The exact host model and power settings
were not preserved, so results describe one recorded environment.

We report median and P95 latency. Median-equivalent rate is
\begin{equation}
R_{\mathrm{eq}}=\frac{1000}{L_{\mathrm{median}}},
\end{equation}
a derived single-record rate rather than measured concurrent throughput. P95
is compared with B1 (10~ms, strict application), B2 (16.7~ms, a 60-record/s
interval), and B3 (100~ms, operational analytics). These study references do
not represent complete sensor-to-application delay or standards compliance.

For $n_f$ transmitted 32-bit numeric fields, modeled payload is
\begin{equation}
P=4n_f \quad \text{bytes per record}.
\end{equation}
Thus, 128, 64, 32, and 16 fields correspond to 512, 256, 128, and 64 bytes.
The estimate excludes framing, timestamps, identifiers, security metadata,
acknowledgments, retransmissions, and packetization.

Within each fold, a 100-tree Random Forest ranks features using only the
fourteen training runs. The same fold-specific ranking is shared by all four
models, which are retrained on the top 64, 32, or 16 features. For the
128-versus-16 comparison, paired fold differences are bootstrapped with
50{,}000 resamples and seed 20260713.

\FloatBarrier
\subsection{Telemetry Degradation and Threshold Transfer}

Random missingness is applied independently to held-out feature cells at
5\%, 10\%, 20\%, and 30\% using seeds 42--46. Group loss sets selected test
features to \texttt{NaN}, after which the training-fitted median imputer is
applied. ``Logs off'' masks the 12 control-panel, relay-log, and Snort
variables; ``PMU off'' masks all PMU/relay measurements. Relay loss masks the
29 measurements associated with R1, R2, R3, or R4. For each fold, the
worst-relay value is the minimum ROC-AUC across those four stress cases.

For central 128-feature models, fixed 0.5 and training-selected thresholds
(maximum F1, maximum Youden's $J$, and training FPR $\leq0.10$) transfer
unchanged to the unseen run. Alert volume is predicted positives per 1{,}000
records.

\FloatBarrier
\section{Results}

\subsection{Random Splits Overstate Unseen-Run Performance}

\begin{table}[H]
\caption{Baseline performance. Random-split values use a 70/30 split;
run-aware values are mean $\pm$ standard deviation across fifteen held-out
runs.}
\label{tab:baseline}
\centering
\scriptsize
\resizebox{\columnwidth}{!}{%
\begin{tabular}{lcccccc}
\toprule
Protocol & Model & Precision & Recall & F1 & FPR & ROC-AUC \\
\midrule
Random & Random Forest & 0.915 & 0.976 & 0.945 & 0.224 & 0.972 \\
Random & XGBoost & 0.871 & 0.950 & 0.908 & 0.346 & 0.922 \\
Random & Feature-space CNN & 0.710 & 1.000 & 0.831 & 1.000 & 0.558 \\
\midrule
Run-aware & Logistic Regression & $0.719\pm0.043$ & $0.976\pm0.005$ &
$0.828\pm0.029$ & $0.934\pm0.020$ & $0.687\pm0.031$ \\
Run-aware & Random Forest & $0.760\pm0.035$ & $0.861\pm0.012$ &
$0.807\pm0.018$ & $0.667\pm0.028$ & $0.677\pm0.029$ \\
Run-aware & XGBoost & $0.763\pm0.038$ & $0.853\pm0.009$ &
$0.805\pm0.022$ & $0.650\pm0.038$ & $0.704\pm0.028$ \\
Run-aware & Compact MLP & $0.743\pm0.037$ & $0.904\pm0.016$ &
$0.815\pm0.024$ & $0.767\pm0.035$ & $0.683\pm0.022$ \\
Run-aware & Feature-space CNN & $0.711\pm0.043$ & $0.998\pm0.006$ &
$0.829\pm0.029$ & $0.998\pm0.004$ & $0.527\pm0.038$ \\
\bottomrule
\end{tabular}
}
\end{table}

Changing the evaluation boundary reduces Random Forest ROC-AUC from 0.972 to
0.677, a 0.295 drop. XGBoost has the highest mean run-aware ROC-AUC among the
four primary models. Logistic Regression and the diagnostic CNN retain
favorable F1 values despite very high FPR, showing that F1 alone can hide weak
unseen-run behavior.

\FloatBarrier
\subsection{Three Models Meet the Strict Processing References}

\begin{table}[H]
\caption{Batch-size-one detector-pipeline latency. Check marks indicate that
P95 falls within the corresponding study reference.}
\label{tab:latency}
\centering
\scriptsize
\resizebox{\columnwidth}{!}{%
\begin{tabular}{lrrrccc}
\toprule
Model & Median & P95 & Med.-eq. & B1 & B2 & B3 \\
 & (ms) & (ms) & rec./s & & & \\
\midrule
Logistic Regression & 1.080 & 1.112 & 926 &
$\checkmark$ & $\checkmark$ & $\checkmark$ \\
Random Forest & 15.669 & 51.711 & 64 &
-- & -- & $\checkmark$ \\
XGBoost & 1.224 & 4.864 & 817 &
$\checkmark$ & $\checkmark$ & $\checkmark$ \\
Compact MLP & 1.160 & 1.690 & 862 &
$\checkmark$ & $\checkmark$ & $\checkmark$ \\
\bottomrule
\end{tabular}
}
\end{table}

Logistic Regression, XGBoost, and the compact MLP meet all three P95
references; Random Forest meets only B3. XGBoost therefore combines the best
mean run-aware discrimination with compliance to the strict study references.

\paperfigure{fig03_latency_envelopes.png}{1.00}
{Batch-size-one detector-pipeline P95 latency for the four primary models,
shown against the 10~ms, 16.7~ms, and 100~ms study references.}
{fig:latency}

\FloatBarrier
\subsection{Sixteen Features Can Preserve Detection While Cutting Payload}

\begin{table}[H]
\caption{Run-aware results by shared training-fold feature budget. Payload
assumes 32-bit numeric fields and excludes protocol overhead.}
\label{tab:feature-budget}
\centering
\scriptsize
\setlength{\tabcolsep}{1.5pt}
\resizebox{\columnwidth}{!}{%
\begin{tabular}{llrrrr}
\toprule
Model & Budget & Features & Payload (B) & ROC-AUC & P95 (ms) \\
\midrule
Logistic Regression & Full & 128 & 512 & 0.687 & 1.112 \\
 & Ranked & 64 & 256 & 0.683 & 1.223 \\
 & Ranked & 32 & 128 & 0.626 & 1.026 \\
 & Ranked & 16 & 64 & 0.629 & 0.781 \\
\midrule
Random Forest & Full & 128 & 512 & 0.677 & 51.711 \\
 & Ranked & 64 & 256 & 0.674 & 28.624 \\
 & Ranked & 32 & 128 & 0.672 & 15.992 \\
 & Ranked & 16 & 64 & 0.721 & 15.353 \\
\midrule
XGBoost & Full & 128 & 512 & 0.702 & 4.864 \\
 & Ranked & 64 & 256 & 0.693 & 0.970 \\
 & Ranked & 32 & 128 & 0.690 & 0.911 \\
 & Ranked & 16 & 64 & 0.708 & 0.818 \\
\midrule
Compact MLP & Full & 128 & 512 & 0.683 & 1.690 \\
 & Ranked & 64 & 256 & 0.683 & 0.899 \\
 & Ranked & 32 & 128 & 0.693 & 0.852 \\
 & Ranked & 16 & 64 & 0.706 & 0.887 \\
\bottomrule
\end{tabular}
}
\end{table}

The 16-feature condition cuts modeled payload eightfold. Its paired
16-minus-128 ROC-AUC differences are $-0.058$ for Logistic Regression,
$+0.045$ for Random Forest, $+0.006$ for XGBoost, and $+0.024$ for the compact
MLP. Bootstrap 95\% intervals are $[-0.071,-0.043]$, $[0.034,0.055]$,
$[-0.002,0.016]$, and $[0.015,0.033]$, respectively. Thus, reduction is
model-dependent. For XGBoost, mean ROC-AUC remains comparable while P95 falls
from 4.864 to 0.818~ms; the small AUC increase is not established by the
interval.

\paperfigure{fig04_payload_latency_quality.png}{1.00}
{Payload--latency trajectories from 128 features to the shared
training-fold-selected 16-feature condition. The figure connects modeled
payload, detector-pipeline P95 latency, and run-aware ROC-AUC; network and
protocol overhead are excluded.}
{fig:tradeoff}

\FloatBarrier
\subsection{Physical Telemetry Dominates Under Degradation}

\begin{table}[H]
\caption{Run-aware ROC-AUC under held-out-test availability degradation.}
\label{tab:availability}
\centering
\scriptsize
\setlength{\tabcolsep}{1.5pt}
\resizebox{\columnwidth}{!}{%
\begin{tabular}{lcccccc}
\toprule
Model & Clean & 10\% Missing & 30\% Missing & Worst Relay & Logs Off & PMU Off \\
\midrule
Logistic Regression & 0.687 & 0.581 & 0.543 & 0.499 & 0.672 & 0.531 \\
Random Forest & 0.677 & 0.680 & 0.659 & 0.580 & 0.674 & 0.506 \\
XGBoost & 0.702 & 0.672 & 0.617 & 0.612 & 0.699 & 0.509 \\
Compact MLP & 0.683 & 0.636 & 0.589 & 0.564 & 0.659 & 0.518 \\
\bottomrule
\end{tabular}
}
\end{table}

Random Forest changes least under random missingness, while XGBoost retains
the strongest fold-wise worst-relay ROC-AUC. Removing cyber/log variables has
little effect, but removing all PMU/relay measurements reduces every model to
near-chance discrimination. The evaluated detectors therefore rely primarily
on physical telemetry.

\paperfigure{fig05_missing_telemetry_robustness.png}{1.00}
{Run-aware F1 under 5\%, 10\%, 20\%, and 30\% test-time random missingness.
Markers show means over fifteen held-out runs and five masking seeds; error
bars show standard deviation across saved fold-seed results.}
{fig:missingness}

\FloatBarrier
\subsection{Central Visibility Helps Until Relay Loss Narrows the Advantage}

\begin{table}[H]
\caption{Model-specific ROC-AUC under analytical edge-visible and central
telemetry conditions.}
\label{tab:placement}
\centering
\scriptsize
\setlength{\tabcolsep}{1.5pt}
\resizebox{\columnwidth}{!}{%
\begin{tabular}{lccccc}
\toprule
Model & Edge R1 & Central clean & Central 10\% miss. & Logs off & Worst relay \\
\midrule
Logistic Regression & 0.531 & 0.687 & 0.581 & 0.672 & 0.499 \\
Random Forest & 0.564 & 0.677 & 0.680 & 0.674 & 0.580 \\
XGBoost & 0.579 & 0.702 & 0.672 & 0.699 & 0.612 \\
Compact MLP & 0.550 & 0.683 & 0.636 & 0.659 & 0.564 \\
\bottomrule
\end{tabular}
}
\end{table}

Central telemetry outperforms the R1 edge-visible condition for every model
when clean, and remains stronger under 10\% random missingness and log loss.
Relay-associated loss narrows this advantage. Logistic Regression falls to
0.499 in its fold-wise worst-relay condition, below its edge-visible result of
0.531; XGBoost retains the strongest degraded-central result at 0.612.

\paperfigure{fig06_edge_vs_central.png}{1.00}
{Run-aware ROC-AUC for the edge-visible R1 condition and central telemetry
under clean and degraded availability. These are analytical feature-visibility
conditions rather than deployed systems.}
{fig:placement}

\FloatBarrier
\subsection{Threshold Transfer Leaves a High Alert Burden}

At the fixed 0.5 threshold, the four-model mean is 856 alerts per 1{,}000
records (FPR 0.755, recall 0.898). Training-based optimization reduces but
does not solve the problem: the FPR-constrained rule still produces 524 alerts
(FPR 0.357) while recall falls to 0.592. Cross-run score separation, not
processing speed alone, remains the limiting factor.

\FloatBarrier
\section{Discussion}

\subsection{Fast Inference Is Necessary, but Not Sufficient}

Three models meet the strict latency references, yet all retain substantial
false-positive behavior on unseen runs. Speed therefore does not establish
operational usefulness. Real-time evaluations should report latency alongside
ROC-AUC, FPR, class balance, and alert volume.

\subsection{Telemetry Reduction Depends on the Model}

Reducing the feature set from 128 to 16 cuts modeled payload from 512 to
64 bytes, but the effect varies by model. Random Forest and the compact MLP
improve, XGBoost remains comparable, and Logistic Regression declines.
Telemetry reduction should therefore be selected within the training boundary
and validated for the intended detector.

\subsection{Central Visibility Increases Channel Dependence}

Central telemetry improves clean-condition discrimination, but relay loss can
narrow or remove this advantage. This trade-off motivates fallback models,
redundant channels, and confidence-aware abstention, although we do not
evaluate these mechanisms here.

\subsection{The Best Trade-Off Is Not Deployment-Ready}

The 16-feature XGBoost configuration provides the strongest measured balance:
64 modeled bytes, 0.708 ROC-AUC, and 0.818~ms P95 latency. However, its
false-positive and threshold-transfer burden remains too high for deployment
without further calibration and system-level testing.

\FloatBarrier
\section{Threats to Validity and Limitations}

The study uses one attack-majority benchmark from one testbed and cannot
establish transfer across utilities, vendors, topologies, or attack taxonomies.
The files lack a valid continuous event sequence, preventing claims about
onset delay, queueing, temporal drift, or early warning.

Latency was measured with cached fold-01 pipelines in one recorded ARM64
macOS environment. The exact host model, CPU-affinity policy, and power mode
were not preserved. Measurements include preprocessing and prediction but
exclude acquisition, serialization, transport, switching, retransmission,
visualization, and operator response. The study references therefore do not
establish IEC~61850 compliance or protection-loop suitability.

Payload counts only assumed 32-bit numeric fields and excludes protocol
overhead. Missingness and relay loss are feature-availability experiments, not
packet-level replay. Feature ranking uses Random Forest impurity importance
and may vary by fold. Edge and central conditions are analytical proxies, and
the fold-wise worst-relay result is a stress case rather than an expected
failure. Threshold results average four models; the diagnostic CNN is excluded
from all communication-budget experiments.

\FloatBarrier
\section{Conclusion}

We examined SCADA intrusion detection as a communication-budgeted service
rather than an accuracy-only classification problem. Run-aware testing exposed
a substantial performance gap hidden by random splitting, while favorable F1
scores sometimes concealed severe false-positive behavior. Three lightweight
models met the strict processing references, and a 16-feature condition
reduced modeled payload eightfold while preserving model-dependent unseen-run
performance. Central telemetry improved clean-condition discrimination, but
its advantage became less reliable under relay-associated feature loss.

These results suggest that lightweight inference is only one part of effective
real-time industrial security. Future systems must jointly support unseen-run
generalization, communication efficiency, placement-aware visibility,
false-alarm control, and resilience to telemetry degradation.


\FloatBarrier
\balance

\begin{thebibliography}{99}

\bibitem{ics-datasets}
T. Morris, ``Industrial Control System (ICS) Cyber Attack Datasets,''
Dataset 1: Power System Datasets, Univ. of Alabama in Huntsville.
Accessed: Jul. 13, 2026. [Online]. Available:
\url{https://sites.google.com/a/uah.edu/tommy-morris-uah/ics-data-sets}

\bibitem{borges2014machine}
R. C. Borges Hink, J. M. Beaver, M. A. Buckner, T. Morris, U. Adhikari,
and S. Pan, ``Machine learning for power system disturbance and cyber-attack
discrimination,'' in \emph{Proc. 7th Int. Symp. Resilient Control Systems
(ISRCS)}, Denver, CO, USA, 2014, Art. no. 6900095,
doi: 10.1109/ISRCS.2014.6900095.

\bibitem{pan2015hybrid}
S. Pan, T. Morris, and U. Adhikari,
``Developing a hybrid intrusion detection system using data mining for power
systems,'' \emph{IEEE Trans. Smart Grid}, vol. 6, no. 6,
pp. 3104--3113, 2015, doi: 10.1109/TSG.2015.2409775.

\bibitem{pan2015disturbances}
S. Pan, T. Morris, and U. Adhikari,
``Classification of disturbances and cyber-attacks in power systems using
heterogeneous time-synchronized data,'' \emph{IEEE Trans. Ind. Informat.},
vol. 11, no. 3, pp. 650--662, 2015,
doi: 10.1109/TII.2015.2420951.

\bibitem{sahani2023machine}
N. Sahani, R. Zhu, J.-H. Cho, and C.-C. Liu,
``Machine learning-based intrusion detection for smart grid computing: A
survey,'' \emph{ACM Trans. Cyber-Phys. Syst.}, vol. 7, no. 2,
pp. 1--31, 2023.

\bibitem{alimi2021review}
O. A. Alimi, K. Ouahada, A. M. Abu-Mahfouz, S. Rimer, and
K. O. A. Alimi, ``A review of research works on supervised learning
algorithms for SCADA intrusion detection and classification,''
\emph{Sustainability}, vol. 13, no. 17, p. 9597, 2021.

\bibitem{upadhyay2021intrusion}
D. Upadhyay, J. Manero, M. Zaman, and S. Sampalli,
``Intrusion detection in SCADA based power grids: Recursive feature
elimination model with majority vote ensemble algorithm,''
\emph{IEEE Trans. Netw. Sci. Eng.}, vol. 8, no. 3,
pp. 2559--2574, 2021.

\bibitem{yang2019deep}
H. Yang, L. Cheng, and M. C. Chuah,
``Deep-learning-based network intrusion detection for SCADA systems,''
in \emph{Proc. IEEE Conf. Commun. Netw. Security}, 2019, pp. 1--7.

\bibitem{wang2022lightweight}
Z. Wang, Z. Li, D. He, and S. Chan,
``A lightweight approach for network intrusion detection in industrial
cyber-physical systems based on knowledge distillation and deep metric
learning,'' \emph{Expert Syst. Appl.}, vol. 206, p. 117671, 2022.

\bibitem{shao2024automated}
J.-M. Shao, G.-Q. Zeng, K.-D. Lu, G.-G. Geng, and J. Weng,
``Automated federated learning for intrusion detection of industrial control
systems based on evolutionary neural architecture search,''
\emph{Comput. Security}, vol. 143, p. 103910, 2024.

\bibitem{arp2022dodo}
D. Arp, E. Quiring, F. Pendlebury, A. Warnecke, F. Pierazzi,
C. Wressnegger, L. Cavallaro, and K. Rieck,
``Dos and don'ts of machine learning in computer security,''
in \emph{Proc. 31st USENIX Security Symp.}, 2022,
pp. 3971--3988.

\bibitem{iec61850-5}
IEC, ``Communication networks and systems for power utility
automation---Part 5: Communication requirements for functions and device
models,'' IEC 61850-5.

\bibitem{iec61850-90-4}
IEC, ``Communication networks and systems for power utility
automation---Part 90-4: Network engineering guidelines,''
IEC TR 61850-90-4.

\bibitem{ieee-c37118}
IEEE, ``IEEE standard for synchrophasor data transfer for power systems,''
IEEE Std C37.118.2-2011, 2011.

\end{thebibliography}

\end{document}
